\section{Limitations}
\label{sec:limitations}
% Keep these. Each one is a real constraint a reviewer will find anyway.

In the disposition family the percentage return is matched across the two
arms of each pair and the dollar amount is not. With the current price and
the share count held fixed, a position 50\% above its basis and one 50\%
below it cannot carry the same unrealised dollar amount. Percentage return is
the quantity the field measurement of the disposition effect is defined on
\citep{odean1998losses}, so it is the quantity matched here, but a model that
responds to dollar magnitude rather than to return would produce an effect in
this family that is not the disposition effect.

The brand salience family cannot claim normative invariance. A real ticker
carries the model's own knowledge of the company, so the two arms of those
pairs differ in information as well as in salience. Both arms use the real
company's sector so that the prompt is not self-contradictory, but the
remaining figures are synthetic in both, which a model with knowledge of the
company may notice. The family is reported separately and never pooled.

Companies in the seven confirmatory families are synthetic. This is
deliberate: \citet{lee2025yourai} report persistent model preferences over
specific real tickers, which would confound every manipulation here. The cost
is that the fundamentals are not real ones, and a model may treat an
unfamiliar name differently from a familiar one in ways the brand salience
family only begins to measure.

Effort is a vendor-defined control, not a measured quantity of reasoning.
\citet{hu2026effort} argue that it sets an upper budget on generation rather
than governing allocation, on a different model family and a different
outcome: they test GPT-OSS at three effort levels on reasoning tasks and
measure the alignment between generation length and human reaction times,
not whether a decision changes. Our token counts agree with their ceiling
claim, so the two results are not in conflict, and we claim an extension of
their manipulation to an outcome they did not measure rather than a
refutation of their finding. We report output token counts per cell so that the
treatment's effect on realised reasoning length is visible, but a null result
on H10 is consistent both with deliberation not reducing bias and with the
effort parameter not delivering the deliberation it nominally buys. These two
are separated only by the token counts, and we are explicit about that.


Synthetic scenarios are not real clients, and a willingness to recommend a
sale in a vignette is not a trade. Effects are prompt level and a different
phrasing of the same cue may move them. Model behaviour changes over time, so
these numbers are tied to the model IDs and run dates in Section~3.4, not to
the model families in general. The human baselines come from field studies of
brokerage accounts, a different instrument from this battery, so the
comparison is indicative rather than like for like. The brand salience and
price path families are exploratory. Twelve pairs per family is the binding
resolution limit of the design: it is what puts study-wide Holm correction out
of arithmetic reach (Section~\ref{sec:results}) and what leaves the
between-family spread in two of six cells indistinguishable from sampling
error (\movedref{sec:rtm}). An extension should add items per family
rather than trials per item. Nothing here is investment advice and
the author is not a licensed adviser.
